\section{From affine two-endpoint bridge to target score}

We consider a general two-endpoint affine bridge of the form
\begin{equation}
x_t = a_t x_0 + b_t x_1 + c_t \varepsilon,
\qquad
\varepsilon \sim \mathcal N(0, I),
\qquad
t \in [0,1],
\end{equation}
with endpoint constraints
\begin{equation}
a_0 = 1,\quad b_0 = 0,\quad c_0 = 0,
\qquad
a_1 = 0,\quad b_1 = 1,\quad c_1 = 0.
\end{equation}
Conditioning on the endpoints $z=(x_0,x_1)$, this defines a conditional Gaussian path
\begin{equation}
p_t(x \mid z) = \mathcal N\!\bigl(x;\mu_t(z), k_t I\bigr),
\end{equation}
where
\begin{equation}
\mu_t(z) = a_t x_0 + b_t x_1,
\qquad
k_t = c_t^2.
\end{equation}
Thus the bridge is simply a Gaussian path with time-dependent mean and isotropic variance.

The quantity of interest is the \emph{time score}
\begin{equation}
s_t(x \mid z) := \partial_t \log p_t(x \mid z).
\end{equation}
Since
\begin{equation}
\log p_t(x \mid z)
=
-\frac{D}{2}\log(2\pi k_t)
-\frac{1}{2k_t}\norm{x-\mu_t(z)}^2,
\end{equation}
we define the standardized residual
\begin{equation}
\epsilon_t(x,z)
:=
\frac{x-\mu_t(z)}{\sqrt{k_t}}.
\end{equation}
A direct differentiation gives
\begin{equation}
\boxed{
\partial_t \log p_t(x \mid z)
=
-\frac{D \dot k_t}{2k_t}
+
\frac{1}{\sqrt{k_t}}\,\dot\mu_t(z)^\top \epsilon_t(x,z)
+
\frac{\dot k_t}{2k_t}\norm{\epsilon_t(x,z)}^2
}
\label{eq:scalar_time_score}
\end{equation}
with
\begin{equation}
\dot \mu_t(z) = \dot a_t x_0 + \dot b_t x_1,
\qquad
\dot k_t = 2 c_t \dot c_t.
\end{equation}
This is the scalar target score used in conditional time-score estimation.

For the vectorized target, note that because the covariance is diagonal, the coordinates are conditionally independent given $z$:
\begin{equation}
p_t(x \mid z)=\prod_{i=1}^D p_t(x_i \mid z).
\end{equation}
Hence
\begin{equation}
\partial_t \log p_t(x \mid z)
=
\sum_{i=1}^D \partial_t \log p_t(x_i \mid z).
\end{equation}
Define the vectorized target score by stacking the coordinatewise terms:
\begin{equation}
\mathrm{vec}\!\bigl(\partial_t \log p_t(x \mid z)\bigr)
:=
\begin{bmatrix}
\partial_t \log p_t(x_1 \mid z) \\
\vdots \\
\partial_t \log p_t(x_D \mid z)
\end{bmatrix}.
\end{equation}
Then the same differentiation as above yields the compact vector formula
\begin{equation}
\boxed{
\mathrm{vec}\!\bigl(\partial_t \log p_t(x \mid z)\bigr)
=
-\frac{\dot k_t}{2k_t}\mathbf 1
+
\frac{1}{\sqrt{k_t}}\,\dot\mu_t(z)\odot \epsilon_t(x,z)
+
\frac{\dot k_t}{2k_t}\,\epsilon_t(x,z)^{\odot 2}
}
\label{eq:vector_time_score}
\end{equation}
where $\odot$ denotes elementwise product and $(\cdot)^{\odot 2}$ denotes elementwise square. Summing all entries of \eqref{eq:vector_time_score} recovers the scalar score in \eqref{eq:scalar_time_score}.

The formulas above immediately cover several useful bridges.

\paragraph{Linear Schr\"odinger-bridge-style path.}
Choose
\begin{equation}
a_t = 1-t,
\qquad
b_t = t,
\qquad
c_t = \sigma \sqrt{t(1-t)}.
\end{equation}
Then
\begin{equation}
\mu_t(z) = (1-t)x_0 + t x_1,
\qquad
k_t = \sigma^2 t(1-t),
\end{equation}
and therefore
\begin{equation}
\dot \mu_t(z) = x_1-x_0,
\qquad
\dot k_t = \sigma^2(1-2t).
\end{equation}
Plugging into \eqref{eq:vector_time_score} gives
\begin{equation}
\mathrm{vec}\!\bigl(\partial_t \log p_t(x \mid z)\bigr)
=
-\frac{1-2t}{2t(1-t)}\mathbf 1
+
\frac{1}{\sigma\sqrt{t(1-t)}}(x_1-x_0)\odot \epsilon_t
+
\frac{1-2t}{2t(1-t)}\epsilon_t^{\odot 2},
\end{equation}
with
\begin{equation}
\epsilon_t
=
\frac{x-[(1-t)x_0+t x_1]}{\sigma\sqrt{t(1-t)}}.
\end{equation}

\paragraph{Cosine-schedule two-endpoint bridge.}
Let
\begin{equation}
s(t)=\frac{1-\cos(\pi t)}{2},
\end{equation}
and define
\begin{equation}
a_t = 1-s(t)=\cos^2\!\left(\frac{\pi t}{2}\right),
\qquad
b_t = s(t)=\sin^2\!\left(\frac{\pi t}{2}\right),
\qquad
c_t = \sigma \sqrt{s(t)(1-s(t))}=\frac{\sigma}{2}\sin(\pi t).
\end{equation}
Then
\begin{equation}
\mu_t(z)=\cos^2\!\left(\frac{\pi t}{2}\right)x_0+\sin^2\!\left(\frac{\pi t}{2}\right)x_1,
\qquad
k_t=\frac{\sigma^2}{4}\sin^2(\pi t),
\end{equation}
and
\begin{equation}
\dot \mu_t(z)=\frac{\pi}{2}\sin(\pi t)(x_1-x_0),
\qquad
\dot k_t=\frac{\pi \sigma^2}{4}\sin(2\pi t).
\end{equation}
Substituting into \eqref{eq:vector_time_score} gives
\begin{equation}
\mathrm{vec}\!\bigl(\partial_t \log p_t(x \mid z)\bigr)
=
-\pi \cot(\pi t)\mathbf 1
+
\frac{\pi}{\sigma}(x_1-x_0)\odot \epsilon_t
+
\pi \cot(\pi t)\,\epsilon_t^{\odot 2},
\end{equation}
where now
\begin{equation}
\epsilon_t
=
\frac{x-\cos^2(\pi t/2)x_0-\sin^2(\pi t/2)x_1}{(\sigma/2)\sin(\pi t)}.
\end{equation}

The general message is simple: once a two-endpoint bridge can be written as a Gaussian affine interpolation
\begin{equation}
x_t = a_t x_0 + b_t x_1 + c_t \varepsilon,
\end{equation}
its scalar and vectorized target scores follow immediately from \eqref{eq:scalar_time_score} and \eqref{eq:vector_time_score}. Different schedules only change $\mu_t$, $k_t$, and their derivatives.